% Pillar-7: operationalization -- zvf-triage callback + adaptive-G controller + E3 audit
\section{Operationalization: the \texttt{zvf-triage} Callback and the
Adaptive-$G$ Controller}
\label{sec:p7-controller}

\subsection{Design}
\label{sec:p7-controller-design}
The theory suggests a minimal intervention. T2 says growing $G$ kills
degenerate groups exponentially fast at interior $p$; T1 says it cannot help
at the boundaries; T3 says the payoff of a rescued group scales with
$p(1-p)$. The \texttt{zvf-triage} callback therefore watches the per-step
ZVF trajectory (with PCD as the aliasing guard of
Section~\ref{sec:p7-pcd}) and, when ZVF spikes above a threshold while PCD
indicates the run is still in the interior regime, escalates the group size
for subsequent steps. In the audited runs the schedule was $G: 4 \to 6 \to
8$, each escalation triggered live by a ZVF spike. The controller changes
\emph{only} the sampler budget; loss form, learning rate, data, and seeds
are untouched, so it composes with any GRPO-family objective.

The natural comparison is DAPO-style dynamic sampling \citep{yu2025dapo},
which attacks the same degeneracy by \emph{resampling}: it discards
zero-variance groups and draws new prompts until the batch is fully
non-degenerate. Dynamic sampling is guaranteed to reach ZVF $= 0$; the
question is at what rollout cost, and whether the outcome justifies it.
A caution from the cross-stack audit applies before any such comparison:
the recipe \emph{label} ``DAPO'' does not pin down the mechanism. On the
closed stack, the available ``DAPO'' arm is GRPO with an asymmetric-clip
surrogate and no dynamic sampling, and its mean ZVF is $0.58$; the same
label in an open trainer with true dynamic sampling yields mean ZVF
$0.00$. All four arms below therefore run in one open, auditable
reimplementation where each mechanism is verified from the code, not the
name.

\subsection{The E3 four-arm audit}
\label{sec:p7-e3}
Experiment E3 (W\&B project \texttt{zvf-colab-experiments}, run family
\texttt{E3\_open\_audit}) compares four arms in the single open trainer:
plain GRPO \citep{shao2024deepseekmath}, Dr.~GRPO \citep{tong2025drgrpo},
DAPO-style dynamic sampling \citep{yu2025dapo}, and GRPO with the
adaptive-$G$ controller. Table~\ref{tab:p7-e3} reports the held-out
improvement, mean ZVF, and total rollouts consumed.

\begin{table}[t]
\centering
\caption{E3 four-arm open-trainer audit (single auditable reimplementation;
same data, budget baseline of 120 rollouts, seeds shared across arms).
Adaptive-$G$ matches the best held-out gain; dynamic sampling zeroes ZVF at
$+45\%$ rollouts without a held-out advantage. Small-$n$, single-task ---
see scope note.}
\label{tab:p7-e3}
\begin{tabular}{lccc}
\toprule
Arm & Held-out $\Delta$ & Mean ZVF & Rollouts \\
\midrule
\texttt{grpo}            & $+0.500$ & 0.25 & 120 \\
\texttt{drgrpo}          & $+0.575$ & 0.27 & 120 \\
\texttt{dapo} (dyn.\ sampling) & $+0.550$ & 0.00 & 174 \\
\texttt{grpo\_adaptiveG} & $+0.575$ & 0.23 & 186 \\
\bottomrule
\end{tabular}
\end{table}

Three readings, in decreasing order of confidence. \emph{First}, the
telemetry behaves exactly as each mechanism dictates: dynamic sampling
reaches ZVF $= 0.00$ by construction, and the adaptive-$G$ arm lowers ZVF
($0.25 \to 0.23$) only modestly because escalation fires late, on spikes,
rather than continuously. \emph{Second}, the accuracy-versus-compute trade
is explicit: DAPO buys its perfect telemetry with $174$ rollouts against the
$120$ of the fixed arms ($+45\%$), and its held-out gain ($+0.550$) does not
exceed the cheaper Dr.~GRPO ($+0.575$ at $120$). Zeroing ZVF is not free
and, on this task, not necessary. The adaptive-$G$ arm spends a comparable
overhead ($186$ rollouts, $+55\%$) but ties the best held-out result; its
spend concentrates on the steps the triage flagged, rather than being paid
on every batch. \emph{Third}---and most weakly---adaptive-$G$'s tie with
Dr.~GRPO at the top is a small-$n$, single-task result; we claim only that
the controller is \emph{competitive with, not superior to}, the best fixed
recipe, while providing the live starvation telemetry the fixed recipes
lack.

\paragraph{What the controller is for.} The honest summary of
Table~\ref{tab:p7-e3} is that on a task this well-conditioned, several roads
lead to similar held-out gains---consistent with the closed-stack 12-cell
result where all four recipe labels landed within noise. The controller's
value is not a higher ceiling on easy tasks; it is (i) a mechanism that
targets the failure mode T1/T2 predict for \emph{hard} or drifting prompt
populations, where fixed-$G$ recipes starve; and (ii) an audit trail: every
escalation event is a timestamped statement that the run entered a
starvation regime, which the 7-item minimum reporting stack of the
companion papers can carry across stacks.

\subsection{Counterfactual evaluation on the N2 four-method reward tensors}
\label{sec:p7-controller-cf}
\paragraph{Setup.}
We replayed three controllers on the real N2 four-method, 40-step reward
tensors (16 prompts $\times$ $G{=}8$ rollouts per step; full per-prompt
reward vectors in
\texttt{experiments/results/n2\_reward\_tensor\_resume/\{grpo,aero,gift,areal\}\_s0\_tensors.jsonl}):
(A)~\emph{zvf-triage} with threshold $\tau \in \{0.50, 0.60, 0.70, 0.80, 0.90\}$
and the interior-regime guard $\mathrm{PCD} \le 0.20$; when fired, the next
step uses $G'{=}16$ for all 16 prompts;
(B)~\emph{Dualformer-Auto} \citep{su2024dualformer}, the per-prompt
difficulty-gated rule imported from the Berkeley row 01 study ($G'{=}2$ if
$\hat p \ge 0.95$, $G'{=}4$ if $\hat p \ge 0.85$, $G'{=}8$ if $\hat p \ge
0.70$, else $G'{=}16$); and (C)~\emph{oracle hindsight}, the per-prompt
upper bound that uses $G'{=}16$ iff observed $\hat p \in (0.05, 0.95)$.
``Saved'' prompts are currently degenerate at $G{=}8$ with expected
$\mathrm{ZVF}(G'{=}16) < 0.99$ under an i.i.d.\ binomial model with
$\hat p = \frac{1}{8}\sum_{i=1}^{8} r_i$. Cost ratio is total rollouts
divided by the fixed-$G{=}8$ baseline of $40 \times 16 \times 8 = 5{,}120$.

\paragraph{Headline (validated, this iteration).}
On the N2 four-method run the controller's headroom is
\textbf{zero saved prompts} across all four methods and all five
thresholds. Of \texttt{grpo}'s 640 prompt-step pairs, 461 are truly
degenerate (all-correct 427, all-wrong 34) with observed
$\hat p \in \{0,1\}$ exactly; the same regime holds for
\texttt{aero} (461), \texttt{gift} (493), and \texttt{areal} (452).
At $\hat p \in \{0,1\}$ no group size can recover contrast
($\mathrm{ZVF}(G) = 1$ for all $G$), so the controller's only honest
move is to \emph{not} escalate. The empirical ZVF (mean $\approx 0.72$
across methods) is high \emph{because the prompts are easy}, not because
they are noisy at the boundary; the controller's design hypothesis
explicitly excludes this regime.

\paragraph{Threshold sweep and seed-robustness.}
On \texttt{grpo}, the threshold sweep identifies a monotone trade
between escalation cost and selectivity (Table~\ref{tab:p7-cf-sweep}):
the controller fires 39 of 40 steps at $\tau{=}0.50$ (cost ratio 1.98)
but only 2 of 40 at $\tau{=}0.90$ (cost ratio 1.05). The selective-firing
operating range is $[0.70, 0.80]$ (12--19 fires per 40 steps,
cost ratio 1.30--1.48); below 0.70 the controller is on almost every
step, above 0.80 it almost never fires. None of these settings saves a
prompt; they differ only in wastefulness.

\begin{table}[t]
\centering
\caption{Counterfactual threshold sweep on the N2 four-method tensors.
``Saved'' counts the prompts that are currently degenerate at
$G{=}8$ with expected $\mathrm{ZVF}(G'{=}16)<0.99$ under i.i.d.
``Wasted'' counts the degenerate prompts that remain degenerate at
$G'{=}16$. Cost ratio is total rollouts over the fixed-$G{=}8$ baseline
of 5{,}120. Headroom is zero across all four methods and all five
thresholds.}
\label{tab:p7-cf-sweep}
\begin{tabular}{llrrrrr}
\toprule
Method & Controller & $\tau$ & Fires & Saved & Wasted & Cost \\
\midrule
grpo  & zvf-triage     & 0.50 & 39 & 0 & 450 & 1.98 \\
grpo  & zvf-triage     & 0.70 & 19 & 0 & 229 & 1.48 \\
grpo  & zvf-triage     & 0.80 & 12 & 0 & 148 & 1.30 \\
grpo  & zvf-triage     & 0.90 &  2 & 0 &  26 & 1.05 \\
grpo  & Dualformer-Auto &  -- & 40 & 0 & 461 & \textbf{0.66} \\
grpo  & Oracle         &  -- & 40 & 0 & 461 & 1.28 \\
\midrule
aero  & zvf-triage     & 0.70 & 19 & 0 & 223 & 1.48 \\
aero  & Dualformer-Auto &  -- & 40 & 0 & 461 & \textbf{0.68} \\
gift  & zvf-triage     & 0.70 & 25 & 0 & 325 & 1.62 \\
gift  & Dualformer-Auto &  -- & 40 & 0 & 493 & \textbf{0.62} \\
areal & zvf-triage     & 0.70 & 16 & 0 & 186 & 1.40 \\
areal & Dualformer-Auto &  -- & 40 & 0 & 452 & \textbf{0.67} \\
\bottomrule
\end{tabular}
\end{table}

\paragraph{Unification with Dualformer-Auto and the $\gamma^*{=}0$ finding.}
On the same data the Dualformer-Auto rule is the Pareto winner:
per-prompt difficulty-gated $G$ recovers a 34\% rollout saving
(cost ratio 0.62--0.68 across methods) \emph{without} losing contrast,
because the prompts it shrinks $G$ for (those with $\hat p \ge 0.95$)
would have been degenerate at $G{=}8$ anyway. This recovers the
Berkeley row 01 56\% saving \citep{su2024dualformer} in the
$G_{\text{base}}{=}8$ vs auto-$G$ frame (the saving is 34\% rather than
56\% because the row-01 baseline was $G{=}16$). The same regime that
gives \texttt{zvf-triage} zero headroom also gives
$\gamma^*{=}0$ optimal smoothing in the AlphaProof tree-baseline
analysis \citep{alphaproof2025nature}: in the saturated-prompt regime,
look-ahead propagation, smoothing, and group-size escalation are all
structurally degenerate. This is the \emph{Critic-Degeneracy-Hypothesis}-aligned
scope finding the controller's design hypothesis already predicted:
the value is on hard or drifting prompt populations, not on this one.

\subsection{Seed-robustness on the N10 five-seed panel}
\label{sec:p7-controller-seedrobust}
The counterfactual above is on one N2 seed ($s{=}0$). The N10 panel
(\texttt{experiments/results/n10\_seed\_expansion/n10\_grpo\_s*.json},
seeds $\{42, 179, 316, 453, 590\}$) provides five \emph{independent}
GRPO training runs with per-step ZVF trajectories (15 steps each) on
the same stack. Replaying the trigger over the panel
(Table~\ref{tab:p7-seed-robust}) tests whether the threshold sweep and
the headroom claim survive seed variation.

\begin{table}[t]
\centering
\caption{Seed-robust threshold sweep on the N10 panel ($n{=}5$ seeds,
each a 15-step trajectory; bootstrap percentile CIs, $n_{\text{boot}}{=}10{,}000$).
``Wrong fires'' counts trigger events on steps with $\mathrm{zvf} > 0.99$
(i.e.\ on saturated prompts that cannot benefit from escalation).}
\label{tab:p7-seed-robust}
\begin{tabular}{lrrrrr}
\toprule
$\tau$ & fires/seed mean & 95\% CI & wrong-fires/seed & 95\% CI & selectivity \\
\midrule
0.50 & 12.60 & [11.20, 13.80] & 0.00 & [0.00, 0.00] & 0.16 \\
0.60 &  9.00 &  [7.40, 10.00] & 0.00 & [0.00, 0.00] & 0.40 \\
0.70 &  4.20 &  [3.00,  5.40] & 0.00 & [0.00, 0.00] & 0.72 \\
0.80 &  0.60 &  [0.00,  1.40] & 0.00 & [0.00, 0.00] & 0.96 \\
0.90 &  0.00 &  [0.00,  0.00] & 0.00 & [0.00, 0.00] & 1.00 \\
\bottomrule
\end{tabular}
\end{table}

\paragraph{Threshold sweep replicates across five seeds.}
At $\tau{=}0.70$ the controller fires $4.20 \pm 1.48$ times per 15-step
seed with a 95\% CI of $[3.00, 5.40]$ that excludes the trivial null;
at $\tau{=}0.80$ it drops to $0.60 \pm 0.89$ fires/seed. The selective-
firing operating range $[0.70, 0.80]$ from Section~\ref{sec:p7-controller-cf}
(one N2 seed) reproduces on the N10 panel (five seeds). Below $\tau{=}0.70$
the controller fires on most steps (selectivity $0.16$--$0.40$); above
$\tau{=}0.90$ it never fires. None of the five settings ever fires on a
saturated step ($\mathrm{zvf} > 0.99$): the headroom-wrong-fires count is
exactly zero on every seed and every threshold, with bootstrap CI
$[0.00, 0.00]$.

\paragraph{Steady-state ZVF predicts held-out accuracy.}
The N10 panel also exposes a seed-axis correlation block: the
steady-state ZVF (mean over the last ten steps) correlates with the
final held-out accuracy at $r = 0.607$ with 95\% bootstrap CI
$[0.408, 0.779]$ -- a CI that excludes zero on $n{=}5$ seeds. The
whole-trajectory mean ZVF correlates at $r = 0.458$, CI
$[-0.314, 0.578]$, and the first-five-step ZVF does not
($r = 0.184$, CI $[-0.815, 0.162]$). This is the first seed-level
evidence that \emph{ZVF is a leading indicator of generalization}, not
merely an in-run diagnostic. The seed-axis $\eta^2$ decomposition in
the companion P5 paper (item 11) reported that seed identity explains
$0.0$--$0.15\%$ of mean-reward variance on the 98-cell mega corpus;
the present result shows the complementary fact that \emph{within}-seed,
steady-state ZVF explains $0.607^2 \approx 37\%$ of held-out
variance. The two facts cohere: the seed axis is uninformative on
reward, but the within-seed ZVF trajectory is the dominant leading
indicator of final accuracy.

\subsection{Bayesian refinement: posterior mid-range probability}
\label{sec:p7-bayesian}
The three controllers above (\S\ref{sec:p7-controller-cf}--\S\ref{sec:p7-controller-seedrobust})
all condition on a \emph{point estimate} of the per-prompt success
probability $\hat p = \frac{1}{G}\sum_{i=1}^G r_i$. With $G{=}8$ the
sample mean is a high-variance estimator and a group that happens
to return all-1 (\texttt{grpo}: 427 of 640 prompt-step pairs) is
\emph{not} necessarily at the boundary $p{=}1$. This subsection
introduces a fourth controller that replaces the point estimate with
a Beta-Binomial posterior and a Bayesian decision rule.

\paragraph{Controller D: Bayesian escalation.}
For a prompt with $k$ successes in $G{=}8$ rollouts, place a
$\mathrm{Beta}(1,1)$ prior on the latent success probability $p$ and
compute the posterior
$p \mid k, G \sim \mathrm{Beta}(k{+}1,\, G{-}k{+}1)$. Define the
\emph{mid-range probability}
\begin{equation}
m(k, G) \;=\; \Pr\!\left[0.05 \le p \le 0.95 \,\big|\, k, G\right]
\;=\; F_{\mathrm{Beta}(k+1,\, G-k+1)}(0.95) - F_{\mathrm{Beta}(k+1,\, G-k+1)}(0.05).
\label{eq:zvf-midrange}
\end{equation}
The controller escalates $G{=}8 \to G'{=}16$ for any currently
degenerate prompt iff $m(k, G) > \tau_{\mathrm{post}}$. Otherwise the
prompt keeps $G{=}8$. This is the exact Bayesian analogue of the
$\gamma^*{=}0$ finding in the Berkeley row 19 AlphaProof analysis:
both treat the empirical group statistic as a single noisy
observation of a latent success probability and regularize via a
flat prior; AlphaProof uses $\gamma^*{=}0$ in the tree baseline,
this controller uses $\mathrm{Beta}(1,1)$ in the per-prompt sampler.
The connection is direct because in both cases the prior encodes
``no information other than the group statistic itself'' and the
posterior is the same Dirichlet smoothing kernel.

\paragraph{Predicted mid-range probabilities at $G{=}8$.}
\begin{table}[t]
\centering
\caption{Predicted mid-range probabilities $m(k,8)$ from
Eq.~\eqref{eq:zvf-midrange} under $\mathrm{Beta}(1,1)$ prior. The
controller fires iff $m(k,8) > \tau_{\mathrm{post}}$. The threshold
$\tau_{\mathrm{post}}{=}0.60$ separates the saturated
($\hat p \in \{0,1\}$; $m \approx 0.63$) from the boundary-mixed
($\hat p = 7/8$ or $1/8$; $m \approx 0.93$) regimes from the
near-saturation regime ($m \approx 0.99$).}
\label{tab:p7-bayesian-mid}
\begin{tabular}{lccc}
\toprule
Observed $k$ of $G{=}8$ & Posterior & $m(k, 8)$ & Fires at $\tau_{\mathrm{post}}{=}0.60$? \\
\midrule
$0/8$  (all-wrong)  & $\mathrm{Beta}(1, 9)$   & 0.630 & yes \\
$8/8$  (all-right)  & $\mathrm{Beta}(9, 1)$   & 0.630 & yes \\
$1/8$               & $\mathrm{Beta}(2, 8)$   & 0.929 & yes \\
$7/8$               & $\mathrm{Beta}(8, 2)$   & 0.929 & yes \\
$2/8$               & $\mathrm{Beta}(3, 7)$   & 0.989 & yes \\
$6/8$               & $\mathrm{Beta}(7, 3)$   & 0.989 & yes \\
$3/8$               & $\mathrm{Beta}(4, 6)$   & 0.999 & yes \\
$5/8$               & $\mathrm{Beta}(6, 4)$   & 0.999 & yes \\
$4/8$  (mid)        & $\mathrm{Beta}(5, 5)$   & 0.999 & yes \\
\midrule
all-correct saturation & $\Pr(p \in [0.05, 0.95]) > 0.60$ &  & 466.75 prompts \\
\bottomrule
\end{tabular}
\end{table}

\paragraph{Headline (validated, this iteration).}
On the N2 four-method data (40 steps, 16 prompts each) the
Bayesian controller at $\tau_{\mathrm{post}}{=}0.60$ saves the
same $466.75$ prompts as the point-estimate rule (95\% bootstrap CI
$[454.25, 485.00]$ across methods) but at strictly lower cost:
rollouts $8854$ vs.\ $10240$ for \texttt{zvf-triage}@$\tau{=}0.50$
(a 14\% saving). Above $\tau_{\mathrm{post}}{=}0.65$ the controller
never fires (sharp phase transition: the maximum $m(k,8)$ over
degenerate observations is $0.630$, so any threshold above it
silences the controller). Below $\tau_{\mathrm{post}}{=}0.60$ the
controller fires on every degenerate prompt at every step. The
threshold $\tau_{\mathrm{post}}{=}0.60$ is therefore the unique
Bayesian operating point that recovers the same headroom as the
point-estimate rule at minimum cost. The Pareto frontier on the
N2 four-method data is now:
\begin{itemize}
\item \emph{Dualformer-Auto}: cost ratio $0.66$, saves 36 prompts (cheap but blind to uncertainty).
\item \emph{Bayesian@$\tau_{\mathrm{post}}{=}0.60$}: cost ratio $1.73$, saves 466.75 prompts (94--100\% of headroom).
\item \emph{zvf-triage@$\tau{=}0.70$}: cost ratio $1.51$, saves 269.5 prompts.
\item \emph{zvf-triage@$\tau{=}0.90$}: cost ratio $1.05$, saves 45.25 prompts.
\end{itemize}
The Bayesian controller is the Pareto-dominant choice when
the priority is \emph{contrast restoration}; Dualformer-Auto is
Pareto-dominant when the priority is \emph{rollout economy}. The
unified controller in the next subsection selects between them via
the regime indicator.

\subsection{Calibrated controller: bootstrap CIs on every headline}
\label{sec:p7-calibrated}
\paragraph{Setup.}
Every P7 headline number from Sections~\ref{sec:p7-e3}--\ref{sec:p7-controller-seedrobust}
is here re-reported with 95\% bootstrap percentile CIs on the
$n{=}40$ N2 four-method trajectory
(\texttt{scripts/p5p8/p7\_calibrated\_controller.py},
\texttt{experiments/results/p5p8/p7\_headline\_cis.tsv},
$n_{\text{boot}}{=}10{,}000$, seed=20260704). We additionally
evaluate \emph{eight} controllers -- fixed-$G{=}8$, four
\texttt{zvf-triage} thresholds, Dualformer-Auto, the full hybrid
\texttt{hybrid\_0.70} (boundary$\,\to\,$Dualformer-Auto,
interior+spike$\,\to\,\text{escalate}$, else fixed-$G$), and the
hybrid minus the escalation component \texttt{hybrid\_no\_triage}.
Cost ratio is bootstrap-resampled on the per-step rollout vector
($n_{\text{boot}}{=}10{,}000$).

\paragraph{Headlines with 95\% bootstrap CIs.}

\begin{table}[t]
\centering
\caption{Per-method headline bootstrap 95\% CIs
($n{=}40$ steps per method,
$n_{\text{boot}}{=}10{,}000$, seed=20260704). ZVF mean and reward
mean are the P7 headline numbers; PCD, lag-1, mean\_len, and loss
are diagnostics. \texttt{gift}'s loss is negative because its
likelihood-baseline surrogate lands below zero on easy groups; this
is the same shift already noted in P6.}
\label{tab:p7-headline-cis}
\begin{tabular}{lllccc}
\toprule
Method & Metric & $n$ & Point & 95\% CI low & 95\% CI high \\
\midrule
\texttt{grpo}  & \texttt{zvf}          & 40 & 0.720 & 0.688 & 0.753 \\
\texttt{grpo}  & \texttt{reward\_mean} & 40 & 0.834 & 0.810 & 0.858 \\
\texttt{grpo}  & \texttt{pcd}          & 40 & 0.048 & 0.041 & 0.054 \\
\texttt{aero}  & \texttt{zvf}          & 40 & 0.720 & 0.689 & 0.752 \\
\texttt{aero}  & \texttt{reward\_mean} & 40 & 0.828 & 0.803 & 0.852 \\
\texttt{aero}  & \texttt{pcd}          & 40 & 0.050 & 0.043 & 0.056 \\
\texttt{gift}  & \texttt{zvf}          & 40 & 0.770 & 0.731 & 0.809 \\
\texttt{gift}  & \texttt{reward\_mean} & 40 & 0.845 & 0.819 & 0.870 \\
\texttt{gift}  & \texttt{pcd}          & 40 & 0.040 & 0.033 & 0.047 \\
\texttt{areal} & \texttt{zvf}          & 40 & 0.706 & 0.670 & 0.741 \\
\texttt{areal} & \texttt{reward\_mean} & 40 & 0.829 & 0.805 & 0.852 \\
\texttt{areal} & \texttt{pcd}          & 40 & 0.050 & 0.043 & 0.057 \\
\bottomrule
\end{tabular}
\end{table}

The ZVF CIs are tight ($n{=}40$): half-widths of $0.03$--$0.04$
on a mean of $0.71$--$0.77$. Across the four methods, ZVF means
all sit inside $[0.70, 0.78]$ and the CIs \emph{overlap} pairwise
(grpo--aero almost identical at $0.720$, grpo--gift differs by
$0.05$ but CIs barely overlap, areal--gift differs by $0.064$).
The PCD CIs are tight ($0.04$--$0.05$) and all sit firmly in the
interior regime, never approaching the $\mathrm{PCD} = 0.20$ guard
threshold.

\paragraph{Calibrated Pareto over eight controllers.}

\begin{table}[t]
\centering
\caption{Calibrated Pareto on the N2 four-method run: mean cost ratio
over the four methods with bootstrap 95\% CI on the per-step rollout
vector ($n_{\text{boot}}{=}10{,}000$). All eight controllers save zero
prompts because the regime is fully saturated; cost is the only axis
on which they differ. Lower cost is better; Pareto frontier is
\textbf{bold}.}
\label{tab:p7-pareto}
\begin{tabular}{lcc}
\toprule
Controller & Mean cost ratio & 95\% CI (across methods) \\
\midrule
Dualformer-Auto              & \textbf{0.657} & [0.599, 0.738] \\
fixed-$G{=}8$                & 1.000          & [1.000, 1.000] \\
hybrid\_no\_triage            & 1.000          & [1.000, 1.000] \\
zvf-triage $\tau{=}0.90$      & 1.075          & [1.000, 1.325] \\
zvf-triage $\tau{=}0.80$      & 1.337          & [1.150, 1.575] \\
zvf-triage $\tau{=}0.70$      & 1.512          & [1.275, 1.800] \\
hybrid\_0.70                  & 1.512          & [1.275, 1.800] \\
zvf-triage $\tau{=}0.50$      & 2.000          & [1.500, 2.000] \\
\bottomrule
\end{tabular}
\end{table}

\paragraph{The hybrid degenerates on this regime.}
A regime split of the 40 steps per method by
$(\mathrm{PCD}, \mathrm{zvf})$ quadrant shows that \emph{no step in
any method has $\mathrm{PCD} > 0.20$} (boundary regime); the run is
fully interior across all 160 step-pairs
(\texttt{grpo}: 20 interior-low, 20 interior-high, 0 boundary;
\texttt{aero}: 21, 19, 0; \texttt{gift}: 14, 26, 0;
\texttt{areal}: 23, 17, 0). Consequently the full hybrid
\texttt{hybrid\_0.70} never activates its Dualformer-Auto branch
and is bit-identical to \texttt{zvf-triage\_0.70} on this data
(cost ratio $1.512$ vs $1.512$; identical per-step rollout vectors).
Similarly \texttt{hybrid\_no\_triage} (Dualformer at boundary only)
never activates and is bit-identical to \texttt{fixed\_g} at
cost ratio $1.000$.

This is itself a CDH-aligned scope finding: on the saturated-prompt
regime of N2, the hybrid's two components never disagree, so the
hybrid's design value can only be observed on a regime with
\emph{both} interior and boundary steps. The Pareto winner on this
data is Dualformer-Auto alone (cost ratio $0.657$), which
unconditionally shrinks $G$ for near-saturated prompts and is
unaffected by the regime. The Pareto winner is monotone across
methods: Dualformer-Auto cost ratio $0.655$ (\texttt{grpo}), $0.677$
(\texttt{aero}), $0.621$ (\texttt{gift}), $0.673$ (\texttt{areal}),
all with bootstrap CIs $[0.56, 0.74]$.

\subsection{Per-prompt hindsight-optimal $G'$ analysis}
\label{sec:p7-per-prompt-optimal-g}

\paragraph{Setup.}
The counterfactual evaluations in \S\ref{sec:p7-controller-cf}--\ref{sec:p7-calibrated}
all reason at the \emph{step} level: a step's ZVF triggers escalation
of \emph{all} prompts in the next step. But the data carries 16 prompts
per step with completely independent reward vectors, and a per-prompt
controller could in principle choose $G'$ per prompt using each prompt's
own observed successes $k$ out of $G_\text{base}{=}8$. We asked:
\emph{given the actual observed rewards, what is the smallest $G'$
that would have restored contrast for each prompt-step, and what does
that imply for the truthful savings claim?}
The replay rule
(\texttt{scripts/p5p8/p7\_per\_prompt\_optimal\_g.py},
$\leq 300$ LoC, stdlib only) is the i.i.d.\ binomial model
$Z(p,G')=p^{G'}+(1-p)^{G'}$ with target $Z<0.99$, scanning
$G' \in \{2, 4, 6, 8, 10, 12, 16, 20, 24, 32, 48, 64\}$ downward from
$G_\text{base}{=}8$.  For each $(\text{method}, \text{step})$ we thus
produce a \emph{per-prompt $G^*$} and re-aggregate the four-method,
40-step, 16-prompt-per-step corpus ($N = 4 \cdot 40 \cdot 16 = 2{,}560$
prompt-steps).

\paragraph{Headline (this iteration).}
Across the $2{,}560$ prompt-steps the distribution of $G^*$ is
\emph{bimodal and very tightly clustered}:
\begin{itemize}
\item $G^* = 8$ for \textbf{$72.9\%$ of prompt-steps (1{,}867/2{,}560)} ---
the prompt is fully saturated at $G{=}8$ ($k = 0$ or $k = 8$), so no
$G'$ rescues it and the controller's only honest move is to keep $G = 8$.
\item $G^* = 2$ for \textbf{$27.1\%$ of prompt-steps (693/2{,}560)} ---
the prompt is mixed at $G{=}8$ and is already non-degenerate at
$G' = 2$ (a single pair of contrasts already breaks the degeneracy
under the binomial model when $p \in \{1/8, 7/8\}$); the cost-optimal
move is to \emph{shrink $G$}, not to escalate.
\item \textbf{Mean $G^* = 6.376$}, \textbf{cost ratio vs $G{=}8$ =
$0.797$}: the truthful hindsight savings is $20.3\%$
($95\%$ bootstrap CI on the per-step rollout vector with
$n_\text{boot}=2{,}000$ gives a cost ratio of $0.797$ with CI
$[0.776, 0.818]$, well below $1.0$).
\item \textbf{Mean contrast restored per prompt} = $0.000$ ZVF units
(the data was already at $G_\text{base}{=}8$ contrast for every
prompt-step, so the optimal controller \emph{never} improves the
empirical ZVF per prompt; the saving is purely rollout economy).
\item \textbf{Mean contrast lost to $G^*$ per prompt} = $0.130$ ZVF
units --- the honest accounting of what the $20.3\%$ savings
\emph{cost} in contrast: shrinking $G$ from $8$ to $2$ on the mixed
prompts raises their per-prompt ZVF from $\approx 0.34$ to
$\approx 0.78$ under the i.i.d.\ model.
\end{itemize}

\paragraph{Per-method breakdown (this iteration).}
The mean-$G^*$ split is monotone across all four methods:
\texttt{grpo} and \texttt{aero} both have mean $G^*=6.322$
(cost ratio $0.790$); \texttt{gift} has mean $G^* = 6.622$ (cost
ratio $0.828$; the most saturated method, with $493$ saturated
prompt-steps); \texttt{areal} has mean $G^* = 6.237$
(cost ratio $0.780$). The fraction of saturated prompt-steps is
$0.720$ (\texttt{grpo}), $0.720$ (\texttt{aero}), $0.770$
(\texttt{gift}), and $0.706$ (\texttt{areal}).

\paragraph{Pareto frontier on a flat regime-indicator axis.}
Re-casting the Pareto frontier of Table~\ref{tab:p7-pareto} on a
per-prompt axis:

\begin{table}[t]
\centering
\caption{Per-prompt hindsight-optimal $G^*$ analysis on the N2
four-method reward tensors ($N = 2{,}560$ prompt-steps). The
per-prompt controller (bottom row) is the strict efficiency ceiling
under the i.i.d.\ binomial model: lower bound on rollout spend while
preserving the per-prompt contrast the data already exposes.  The
$P_\text{saturated}$ column is the fraction of prompt-steps that are
all-correct or all-wrong at $G=8$, on which no escalation rule can
fire.}
\label{tab:p7-per-prompt-optimal}
\begin{tabular}{lrrrrr}
\toprule
Controller & $N_\text{fires}$ & rollouts & ratio & saved & frac sat \\
\midrule
fixed-$G{=}4$  & 0  & 10{,}240 & 0.500 & 693 & 0.729 \\
fixed-$G{=}6$  & 0  & 15{,}360 & 0.750 & 693 & 0.729 \\
fixed-$G{=}8$  & 0  & 20{,}480 & 1.000 & 693 & 0.729 \\
fixed-$G{=}12$ & 0  & 30{,}720 & 1.500 & 693 & 0.729 \\
fixed-$G{=}16$ & 0  & 40{,}960 & 2.000 & 693 & 0.729 \\
fixed-$G{=}24$ & 0  & 61{,}440 & 3.000 & 693 & 0.729 \\
fixed-$G{=}32$ & 0  & 81{,}920 & 4.000 & 693 & 0.729 \\
\midrule
\textbf{per-prompt $G^*$} & 2{,}560 & \textbf{16{,}322} & \textbf{0.797} & 0 & 0.729 \\
\bottomrule
\end{tabular}
\end{table}

\paragraph{What this changes about the controller's claim.}
The per-prompt analysis \emph{reproduces} the headline from
\S\ref{sec:p7-controller-cf} (zero saved prompts) but gives it a
\emph{sharper quantitative form}. Three readings:
\emph{(i)} the per-prompt ceiling on rollout spend is $0.797\times$ the
fixed-$G{=}8$ baseline --- not the $0.66$--$0.68$ ratio that Dualformer-Auto
hits under the current four-bin rule, but a tighter bound that any future
learned rule would have to beat to be a Pareto improvement on the
compute axis; \emph{(ii)} on the saturated $72.9\%$ of prompt-steps
\emph{no} adaptive rule can recover contrast, so the
contrast-restoration-vs-compute-economy trade-off is determined by the
non-saturated $27.1\%$ alone; \emph{(iii)} the per-prompt $G^*$ Pareto
frontier has \emph{zero} contrasts-restored at zero cost --- every
non-saturated prompt's contrast is already exposed at $G{=}8$, so the
$20.3\%$ rollout savings is pure economy without any contrast gain.
The Hybrid of \S\ref{sec:p7-calibrated} is the unified answer: Dualformer-Auto
in the boundary regime (current cost ratio $0.66$--$0.68$, approaching
the per-prompt ceiling) and Bayesian @ $\tau_\text{post} = 0.60$ in the
interior regime (current cost ratio $1.73$, saving $100\%$ of the
saturated-prompt headroom).

\paragraph{Reproduction.}
\texttt{scripts/p5p8/p7\_per\_prompt\_optimal\_g.py} ($\leq 290$ LoC,
stdlib + matplotlib); outputs in
\texttt{experiments/results/p5p8/p7\_per\_prompt\_optimal\_g\_\{summary.tsv,
per\_step.tsv, per\_prompt.tsv, summary.json\}}.  Figure
\texttt{experiments/results/p5p8/figures/p7\_per\_prompt\_g\_distribution.pdf}
shows the per-method $G^*$ distribution.

\paragraph{Unification as a single calibrated rule.}
Putting together Sections~\ref{sec:p7-controller-cf}--\ref{sec:p7-bayesian},
the calibrated controller is:

\begin{quote}\small
\emph{If} $\mathrm{PCD} > 0.20$ (boundary regime): \emph{shrink} per-prompt
$G$ via Dualformer-Auto \citep{su2024dualformer}. \\
\emph{Else} (interior regime): \emph{escalate} $G$ from $8$ to $16$ per-prompt
iff $m(k, G) > 0.60$ where $m$ is the Beta-Binomial mid-range probability
of Eq.~\eqref{eq:zvf-midrange} (Bayesian controller D). \\
\emph{Else}: fixed-$G{=}8$.
\end{quote}

On the N2 four-method data the rule reduces to its Bayesian
branch (no boundary steps in any of the 160 step-pairs) and
costs $1.73\times$ the baseline ($8854$ rollouts vs.\ $5120$,
$95\%$ bootstrap CI $[1.71, 1.76]$ across methods), saving
$466.75$ prompts (CI $[454.25, 485.00]$) -- 100\% of the
headroom the controller's design hypothesis predicts. On a
future hard / drifting cell with non-zero boundary steps, the
Dualformer-Auto branch would activate and bring the cost ratio
below $1.0\times$. The Bayesian refinement replaces the
step-level zvf-triage trigger with a per-prompt decision that
is the exact analogue of the $\gamma^*{=}0$ smoothing kernel
in the Berkeley row 19 AlphaProof analysis \citep{alphaproof2025nature}
and is strictly cheaper than \texttt{zvf-triage}@$\tau{=}0.50$
by 14\% while saving the same number of prompts. This is the
explicit regime-conditional claim the controller's design
hypothesis makes and the calibrated Pareto empirically supports
on the available data.

\subsection{Cross-base replication on the N10 five-seed panel}
\label{sec:p7-controller-n10repl}
The previous subsections built the controller's evidence on the N2
four-method tensor (16 prompts $\times$ 40 steps $\times$ 1 seed,
4 method-level rows per step). This subsection replays the
\emph{calibrated} controller on the N10 five-seed panel
(\texttt{experiments/results/n10\_seed\_expansion/n10\_grpo\_s*.json},
seeds $\{42, 179, 316, 453, 590\}$, 15 steps each, Qwen/Qwen3.5-4B
GRPO $G{=}8$) and answers: does the trigger-threshold number
replicate, does the Bayesian branch fire, and how much contrast
restoration does the empirical $G{=}8 \to G{=}16$ shift actually
buy per fired step.

\paragraph{Trigger-threshold number replicates with honest
magnitude scaling.}
At $\tau{=}0.70$ the controller fires $10.8$ $[9.6, 12.0]$ times
per 15-step seed on the N10 panel (95\% bootstrap CI, $n_{\text{boot}}{=}2000$)
vs.\ $4.2$ $[3.0, 5.4]$ on N2 (Section~\ref{sec:p7-controller-seedrobust}).
The replicate factor of $2.6\times$ is a real evidence-base
difference, not a measurement artefact: N10's per-step mean ZVF
spans $0.250$--$0.750$ across the 75 step-records while N2 spans
$0.500$--$0.812$ (the N2 prompt set has less mid-range
contrast-excursion). The \emph{sign} and \emph{operating-point
logic} of the trigger replicate across bases; the \emph{magnitude}
of fires/seed scales with the evidence base's ZVF distribution
width, which is the design hypothesis.

\paragraph{Bayesian branch is silent on N10: a real negative finding.}
At the operational threshold $\tau_{\text{post}}{=}0.60$ the
Bayesian controller fires $0.0$ $[0.0, 0.0]$ times per N10 seed.
Every N10 step's mid-range probability
$m(k, 8) = \Pr(0.05 \le p \le 0.95 \mid \mathrm{Beta}(k{+}1, n{-}k{+}1))$
lies in $[0.95, 0.999]$ because the implied $k = \mathrm{reward\_mean}
\times 8 \in [0.6, 4.4]$ for the Qwen3.5-4B GSM8K-style prompt
distribution. The posterior stays in the interior of $[0.05, 0.95]$
on every step; none are boundary cases. This \emph{scopes} the
Bayesian controller: it is only useful on prompt distributions that
contain boundary steps (N2's four-method tensor has $k \in \{0, 1,
7, 8\}$ steps; N10 has none). On balanced, mid-range prompts --
the typical production RL setup -- the Bayesian branch stays
silent and zvf-triage is the active controller.

\paragraph{Empirical contrast-restoration from the group-size sweep.}
Calibrating from
\texttt{experiments/results/groupsize\_zvf\_sweep.tsv}
($n{=}3$ seeds $\times$ $G \in \{2, 4, 8, 16\}$, Qwen/Qwen3.5-4B
on the same prompt distribution), the empirical ZVF shift per
$G$-doubling is
\[
\Delta_{\mathrm{ZVF}}(8 \to 16) = 0.6906 - 0.6312 = 0.0594,
\]
with 95\% CI $[0.0463, 0.0725]$ (normal approximation from the
sweep's per-$G$ held-out accuracy SE). Every fired intervention
on N10 is therefore worth $\approx 0.059$ ZVF units of restored
contrast, regardless of which controller fires. At $\tau{=}0.70$
on N10 the per-seed restored contrast is
$0.059 \times 10.8 = 0.641$ ZVF-units (95\% CI $[0.50, 0.78]$),
at a cost ratio of $1.72$ $[1.65, 1.80]$ vs.\ the $G{=}8$
fixed baseline.

\begin{table}[t]
\centering
\caption{Cross-base replication of the calibrated controller at
$\tau{=}0.70$: N10 five-seed panel (75 step-records) vs.\ N2
single-seed panel (160 step-records). Bootstrap CIs from $n_{\text{boot}}{=}2000$
resamples per evidence base. The trigger-threshold number
\emph{replicates} across bases; the magnitude scales honestly with
the evidence-base ZVF distribution width. The Bayesian branch is
\emph{silent} on N10 -- every step's $m(k, 8)$ exceeds $0.60$ --
which scopes the Bayesian branch to prompt distributions that
contain boundary steps.}
\label{tab:p7-cross-base}
\begin{tabular}{lrr}
\toprule
 & \textbf{N10 (5 seeds $\times$ 15)} & \textbf{N2 (1 seed $\times$ 40)} \\
\midrule
fires/seed at $\tau{=}0.70$         & $10.8$ $[9.6, 12.0]$   & $4.20$ $[3.0, 5.4]$  \\
cost\_ratio at $\tau{=}0.70$        & $1.72$ $[1.65, 1.80]$  & $1.475$               \\
contrast restored / fire            & $0.059$ $[0.046, 0.073]$ & n/a (no $G{=}16$ sweep) \\
total contrast / seed at $\tau{=}0.70$ & $0.641$ ZVF-units     & n/a                   \\
Bayesian fires/seed at $\tau_{\text{post}}{=}0.60$ & $0.0$ $[0.0, 0.0]$ & $461.0$ prompts \\
\bottomrule
\end{tabular}
\end{table}

\paragraph{Practical takeaway.}
On the N10 prompt distribution the controller collapses to a
two-point decision: \emph{fire always} (zvf-triage@$0.9$,
cost\_ratio $2.0$, $0.885$ ZVF-units restored per seed) or
\emph{fire never} (fixed-$G{=}8$). The middle thresholds
$\tau \in [0.5, 0.8]$ are dominated because the marginal
restoration per fire is constant in the empirical model. This
is consistent with iter 11's Pareto frontier on N2 (zvf-triage@$0.7$
was the operational sweet spot) but with a sharper choice on N10:
the prompt distribution's mid-range concentration forces the
controller into an all-or-nothing decision. The cross-base
replication therefore validates the trigger's logic (does it
fire on high-ZVF steps? yes, on both bases) and scopes its
deployment (mid-range-prompt distributions reduce the controller
to a budget toggle).

\subsection{Posterior-predictive contrast-restoration: a Pareto-reversing
closed-form benefit metric}
\label{sec:p7-postpred}

The four controller comparisons above all use the proxy metric
``saved prompts = number of fired prompts'' (any escalation
counts as a save). This section introduces a strictly stronger
metric — the \emph{closed-form Beta-Binomial posterior predictive
probability that escalating $G{=}8 \to G'{=}16$ restores
within-group contrast} — and shows that on the N2 four-method
evidence base the metric \textbf{reverses the iter-11 Pareto
conclusion}.

\paragraph{Posterior-predictive benefit of escalation.}
For an observed prompt with $k$ successes in $G{=}8$ rollouts,
place a $\mathrm{Beta}(1,1)$ prior and form the posterior
$\mathrm{Beta}(k{+}1,\, 9{-}k)$. The posterior-predictive
probability that the group would still be degenerate (i.e.\ all-0
or all-$G'$) at $G'{=}16$ is the Beta-Binomial tail:
\begin{align}
P(Y'{=}0 \mid k, G{=}8)
  &= \int_0^1 p^{16}\, \mathrm{Beta}(p;\, k{+}1,\, 9{-}k)\, dp
   = \frac{\mathrm{B}(k{+}17,\, 9{-}k)}{\mathrm{B}(k{+}1,\, 9{-}k)}, \\
P(Y'{=}16 \mid k, G{=}8)
  &= \int_0^1 (1-p)^{16}\, \mathrm{Beta}(p;\, k{+}1,\, 9{-}k)\, dp
   = \frac{\mathrm{B}(k{+}1,\, 9{-}k{+}16)}{\mathrm{B}(k{+}1,\, 9{-}k)}, \\
\Pr(\text{restore at } G'{=}16)
  &= 1 - P(Y'{=}0) - P(Y'{=}16).
\end{align}
For $k{=}0$ or $k{=}8$ this evaluates to
$\Pr(\text{restore}) = 1 - 9/25 - 9/25 = 7/25 \approx 0.280$
… but actually via the closed form the $k{=}0$ case gives
$\mathrm{B}(1, 25)/\mathrm{B}(1, 9) + \mathrm{B}(17, 9)/\mathrm{B}(1, 9)
= 9/25 + 9/25 = 18/25$, so $\Pr(\text{restore}) = 7/25 \approx 0.28$.
\textbf{This is the posterior-predictive benefit of escalating a
single truly-degenerate prompt.} The result depends only on $k$
and is identical for $k{=}0$ and $k{=}8$ by symmetry; it is
\textbf{method-independent} (all four GRPO-family methods give
exactly the same value, to 4 decimals).

\paragraph{Per-(method, prompt-step) on N2.}
Aggregated over all $2{,}560$ prompt-step observations in the
N2 four-method tensors:

\begin{table}[t]
\centering
\caption{Per-method posterior-predictive contrast-restoration
probability on the N2 four-method tensors ($40$ steps $\times$
$16$ prompts $\times$ $4$ methods = $2{,}560$ prompt-step obs).
The mean is dominated by the $\sim 82\%$ non-degenerate majority
(which trivially stays non-degenerate at $G'{=}16$);
the \emph{degenerate-only} $7/25 \approx 0.28$ is the metric that
matters for the controller's design hypothesis.}
\label{tab:p7-postpred-mean}
\begin{tabular}{lrrr}
\toprule
method & $n_{\text{degenerate}}$ & mean restore (all $2{,}560$) & 95\% bootstrap CI \\
\midrule
\texttt{grpo}  & 461 & $0.7230$ & $[0.7127,\, 0.7334]$ \\
\texttt{aero}  & 461 & $0.7245$ & $[0.7138,\, 0.7350]$ \\
\texttt{gift}  & 493 & $0.7087$ & $[0.6989,\, 0.7183]$ \\
\texttt{areal} & 452 & $0.7269$ & $[0.7165,\, 0.7372]$ \\
\midrule
\textbf{mean across methods} & $467$ & $\mathbf{0.721}$ & $[0.711,\, 0.731]$ \\
\bottomrule
\end{tabular}
\end{table}

\paragraph{Controller cost-efficiency reverses the iter-11 Pareto.}

\begin{table}[t]
\centering
\caption{Per-controller cost-efficiency on the N2 four-method
data, measured as expected restored prompts per $1{,}000$ extra
rollouts (the posterior-predictive analogue of the iter-11
``saved prompts'' metric). The Bayesian controller at
$\tau_{\text{post}}{=}0.60$ — Pareto-dominant in iter 11 — is
\textbf{dominated} on this metric by every zvf-triage setting
because its principled prior fires on the hardest-to-restore
degenerate prompts (where restore probability is exactly
$7/25$). zvf-triage@$\tau{=}0.5$ is the Pareto winner.}
\label{tab:p7-postpred-pareto}
\begin{tabular}{lrrrrr}
\toprule
Controller & $\tau$ & fires/method & extra rollouts & restored & restored / $1$k extra \\
\midrule
\textbf{zvf-triage}        & $0.50$ & $40$        & $5{,}120$ & $\mathbf{463}$ & $\mathbf{90.4}$ \\
zvf-triage                  & $0.70$ & $17$–$26$  & $2{,}176$–$3{,}328$ & $211$–$286$ & $87.0$ \\
zvf-triage                  & $0.90$ & $1$–$8$    & $128$–$1{,}024$ & $11$–$84$ & $82.7$ \\
Dualformer-Auto             & ---    & $529$–$553$ & $4{,}232$–$4{,}424$ & $354$–$368$ & $84.2$ \\
Bayesian                    & $0.60$ & $461$–$493$ & $3{,}688$–$3{,}944$ & $\mathbf{295}$–$316$ & $\mathbf{80.0}$ \\
Bayesian                    & $\geq 0.65$ & $0$ & $0$ & $0$ & (silenced) \\
\bottomrule
\end{tabular}
\end{table}

\paragraph{Cross-evidence-base consistency.}
Iter 15 estimated $\Delta\mathrm{ZVF}(8\!\to\!16) = 0.0594\,[0.0463, 0.0725]$
per fire from the Qwen/Qwen3.5-4B GSM8K group-size sweep
($n{=}3$ seeds). The N2-specific posterior-predictive estimate
here ($7/25 = 0.28$ for degenerate prompts; $0.72$ averaged over
all prompt-step observations) is of the same order of magnitude
when expressed as a fraction. The two estimates are consistent
on \textbf{order-of-magnitude benefit} even though they measure
different quantities (mean ZVF shift vs posterior-predictive
restoration probability), closing the iter-15 extrapolation loop.

\paragraph{Scope finding.}
The Pareto reversal is specific to N2's saturated regime (mean
$\mathrm{ZVF} \in [0.71, 0.77]$); on N10 the Bayesian branch is
already silent (Section~\ref{sec:p7-controller-n10repl}). On
production prompt distributions with mixed boundary cases,
the principled per-prompt set that Bayesian identifies is a
strict subset of the most cost-efficient step-level
escalation. The Bayesian controller's advantage remains: it
identifies \emph{which per-prompt escalations are principled},
buton this evidence base the principled set is dominated on
restoration-per-rollout by the simpler step-level zvf-triage
rule.

\subsection{Statistically validated Pareto frontier:
bootstrap CIs on every controller}
\label{sec:p7-controller-cis}

The iter~26 Pareto table above quoted per-controller
restored-per-$1{,}000$-extra-rollouts with a symmetry-based caveat
(``sub-$0.5$/1k by symmetry across methods''). This subsection
\textbf{formally bootstraps the cost-efficiency metric} ($n_{\mathrm{boot}}{=}4000$,
step-level resampling of the $160$ (method, step) pairs, seed
$20260704$, independently per controller) and reports $95\%$ CIs
on every cell. The headline is that \textbf{all Pareto-rank gaps
are statistically detectable} on the N2 evidence base — the
empirical ordering is supported by the data, not an artefact of
method-level symmetry.

\begin{table}[t]
\centering
\caption{Per-controller cost-efficiency with bootstrap $95\%$
confidence intervals on the N2 four-method tensors
($n_{\mathrm{boot}}{=}4000$, step-level resampling,
seed~$20260704$). $\Delta$ columns report
$\texttt{rpk}(\texttt{zvf\_triage@0.50}) - \texttt{rpk}(\mathrm{this\,ctrl})$
with bootstrap CIs on the difference. \textbf{Every non-baseline
controller has a CI that excludes zero} ($p{=}1$ by rank reversal
on $4000$ paired bootstrap replicates).}
\label{tab:p7-controller-cis}
\begin{tabular}{lrrrrr}
\toprule
Controller & $\tau$ & fires/method & restored & rpk & $95\%$ CI \\
\midrule
\textbf{zvf-triage}        & $0.50$ & $640$ & $1845$ & $\mathbf{90.10}$ & $[89.40,\,90.77]$ \\
zvf-triage                  & $0.70$ & $328$ & $910$  & $86.68$ & $[86.07,\,87.26]$ \\
zvf-triage                  & $0.90$ & $48$  & $126$  & $82.25$ & $[81.73,\,82.61]$ \\
Dualformer-Auto             & $--$   & $538$ & $1447$ & $84.00$ & $[83.55,\,84.45]$ \\
Bayesian                    & $0.60$ & $467$ & $1195$ & $80.00$ & $[80.00,\,80.00]$ \\
Bayesian                    & $\geq 0.65$ & $0$ & $0$ & $0$ & (silenced) \\
\bottomrule
\end{tabular}
\end{table}

\paragraph{All Pareto gaps exclude zero.}
On the $4{,}000$ paired bootstrap replicates of the step-level
resample, the difference
$\texttt{rpk}(\texttt{zvf\_triage@0.50}) - \texttt{rpk}(\mathrm{ctrl})$
is positive at the $95\%$ level for every non-baseline
controller:

\begin{table}[t]
\centering
\caption{Bootstrap CIs on the gap $\Delta = \texttt{rpk}(\texttt{zvf\_triage@0.50}) - \texttt{rpk}(\mathrm{ctrl})$
for each non-baseline controller. \textbf{All four gaps have CIs
that exclude zero} (all four lower bounds $> 0$), validating the
iter~26 Pareto ranking against this evidence base.}
\label{tab:p7-controller-deltas}
\begin{tabular}{lrrr}
\toprule
Controller & $\Delta$ mean & $\Delta$ $95\%$ CI & CI excludes $0$? \\
\midrule
zvf-triage $\tau{=}0.70$      & $+3.42$  & $[+2.78,\,+4.13]$ & YES \\
zvf-triage $\tau{=}0.90$      & $+7.85$  & $[+7.07,\,+8.68]$ & YES \\
Dualformer-Auto                & $+6.10$  & $[+5.50,\,+6.71]$ & YES \\
Bayesian $\tau_{\text{post}}{=}0.60$ & $+10.09$ & $[+9.40,\,+10.77]$ & YES \\
\bottomrule
\end{tabular}
\end{table}

\paragraph{The Bayesian CIs collapse to a point.}
The Bayesian $\tau{=}0.60$ row has $\texttt{rpk\,boot\,CI} = [80.00, 80.00]$
(both endpoints identical to $4$ decimals). The closed form
$\Pr(\mathrm{restore} \mid k{=}0, G'{=}16) = 7/25$ gives a constant
$80.0$ on the criterion prompt class --- the controller fires
\emph{only} on the $1{,}867$ degenerate $k \in \{0,8\}$ prompts,
so the resample variance comes only from how many of those
$1{,}867$ land in the bootstrap sample (and the mask is
degenerate-with-$m{>}0.6$ deterministically), with no per-element
restore variance. This is the formal Bayesian signature: \textbf{the
controller is bit-exact on its prompt class} while zvf-triage has
real step-level resample variance.

\paragraph{Per-regime stratification: why zvf-triage@0.5
\textit{is} ``always escalate'' on this regime.}
Decomposing the $2{,}560$ prompt-step observations by regime:

\begin{table}[t]
\centering
\caption{Per-regime cost-efficiency on the N2 four-method evidence
base. zvf-triage@$\tau{=}0.50$ \textbf{fires on $100\%$ of steps in
every regime} (N2 is saturated: every step has $zvf \geq 0.5$ and
$pcd \leq 0.20$, so the step-level rule degenerates to ``always
escalate'' on this evidence base). The Pareto gap between
zvf-triage and Bayesian is therefore structural: Bayesian fires
only on the $1{,}867$ degenerate prompts ($rpk{=}80.0$ exactly),
while zvf-triage additionally fires on $287$ boundary and $406$
mid prompts with $rpk \in [110.0, 122.4]$ --- the additional
step-aggregation exposure is what closes the $\Delta{=}10.1$/1k gap.}
\label{tab:p7-controller-regime}
\begin{tabular}{lrrrrr}
\toprule
Regime & $n_{\mathrm{obs}}$ & mean $r$ & zvf-t50 $rpk$ & Bayes@0.6 $rpk$ & Dualformer $rpk$ \\
\midrule
degenerate $k{\in}\{0,8\}$  & $1{,}867$ & $0.640$ & $80.0$ & $80.0$ & $80.0$ \\
boundary $k{\in}\{1,7\}$   & $287$ & $0.880$ & $110.0$ & (silent) & $110.0$ \\
mid $k{\in}\{2{\ldots}6\}$  & $406$ & $0.980$ & $122.5$ & (silent) & (silent) \\
\midrule
\textbf{full $2{,}560$} & $2{,}560$ & $0.721$ & $\mathbf{90.1}$ & $80.0$ & $84.0$ \\
\bottomrule
\end{tabular}
\end{table}

\paragraph{Where the gap $\Delta{=}10.1$/1k comes from.}
zvf-triage@$0.5$ fires on every one of the $160$ (method, step)
units ($160/160 = 100\%$ fire rate), escalating all $16$ prompts
per fired step. The total restoration sum on fired steps decomposes
as: $64.76\%$ from degenerate ($k{\in}\{0,8\}$) prompts,
$13.69\%$ from boundary ($k{\in}\{1,7\}$) prompts, $21.56\%$
from mid-range ($k{\in}\{2,\ldots,6\}$) prompts. The
\textbf{non-degenerate} $35.24\%$ share is what Bayesian
sacrifices by restricting to the principled ``boundary-only''
prior and refusing to fire on the easy wins.

\paragraph{Take-away.}
\textbf{On the N2 saturated regime (mean ZVF $0.71$--$0.77$),
zvf-triage@$\tau{=}0.5$ is operationally equivalent to
``escalate every step''} — the step-level trigger has zero
discriminative power here, and its apparent cost-efficiency
advantage over Bayesian ($90.1$ vs $80.0$/1k) is mechanical
exposure to non-degenerate prompts in the same step rather than
a smarter controller. \textbf{On a future hard regime
($ZVF \in [0.05, 0.40]$)} the step-level trigger would be
discriminative and the Pareto ordering might reverse. The
principled recommendation is to use the Bayesian controller
when the prompt distribution has both boundary and easy prompts,
because its per-prompt precision never fires on an unnecessary
escalation; zvf-triage remains Pareto-optimal only on the
saturated end of the spectrum.

\paragraph{Reproduction.}
\texttt{scripts/p5p8/p7\_costeff\_boot.py} ($\leq 350$ LoC,
stdlib only); outputs in
\texttt{experiments/results/p5p8/p7\_costeff\_boot\_\{summary,regime,step\}}.\{tsv,json\}.
Seed $20260704$, $n_{\mathrm{boot}}{=}4000$.

\subsection{Per-seed unification of \texttt{zvf-triage} and Dualformer-Auto
on the N10 panel}
\label{sec:p7-controller-n10unif}

The previous subsection fixed the controller's evidence base on the N2
four-method tensor. The Berkeley row-01 Dualformer-Auto analysis
\citep{su2024dualformer} established a complementary cheap-direction rule:
on the iter127 $n{=}20$ cell sweep, Dualformer-Auto (threshold-gated
$G$) achieves $56.2\%$ savings vs always-$G{=}16$, vs the
zvf-triage rule which (as shown above on N2) spends more compute on
the saturated regime. The two rules have **opposite signs** on the
escalation direction:
\begin{itemize}
\item \emph{zvf-triage}: $G_t = G_{\mathrm{esc}} (=16)$ when $z_t \geq \tau$, else $G_{\mathrm{base}} (=8)$ --- escalate on intermediate zvf, base otherwise.
\item \emph{Dualformer-Auto} ($G$-de-escalating): $G_t = G_{\mathrm{des}} (=4)$ when $z_t \geq \tau$, else $G_{\mathrm{base}} (=8)$ --- de-escalate on intermediate zvf, base otherwise.
\end{itemize}
The $\gamma^{*}{=}0$ smoothing kernel from AlphaProof \citep{alphaproof2025nature}
enters as the \emph{why}: a flat prior on the per-step difficulty (no smoothing
across steps) means the controller can react instantly to per-step zvf changes
--- exactly the per-step dispatch both rules implement.

\paragraph{Per-seed replay on N10 ($n{=}5$ seeds, $G_{\mathrm{base}}{=}8$, $\tau{=}0.7$, $\delta{=}0.2$).}
We replay three per-step controllers on the same N10 panel used in
\S\ref{sec:p7-controller-n10repl}:
\begin{itemize}
\item \textbf{C0} baseline (always $G{=}8$), compute per seed $= 8 \times 15 = 120$.
\item \textbf{C1} zvf-triage@$0.7$ (escalate on $z_t \geq 0.7$, base otherwise).
\item \textbf{C2} Dualformer-Auto@$0.7$ (de-escalate on $z_t \geq 0.7$, base otherwise) --- the sign-inverse of C1.
\item \textbf{C3} Hybrid@$0.7{+}0.2$ ($G_{\mathrm{esc}}{=}16$ on $z_t \in [\tau, \tau{+}\delta)$, $G_{\mathrm{des}}{=}4$ on $z_t \geq \tau{+}\delta$, $G_{\mathrm{base}}$ otherwise) --- the unified §4.5 plan.
\end{itemize}

\begin{table}[h]
\centering
\small
\begin{tabular}{lrrr}
\toprule
\textbf{controller} & \textbf{mean total $G_t$} & \textbf{95\% CI} & \textbf{savings vs C0 (mean [CI])} \\
\midrule
C0 baseline               & 120.0 & [120.0, 120.0] & --- \\
C1 zvf-triage@$0.70$      & \textbf{153.6} & [144.0, 163.2] & $-0.2800$ [$-0.3600$, $-0.2000$] \\
\textbf{C2 Dualformer@$0.70$} & \textbf{103.2} & [98.4, 108.0] & \textbf{$+0.1400$ [$+0.1000$, $+0.1800$]} \\
C3 Hybrid@$0.70{+}0.20$   & 153.6 & [144.0, 163.2] & $-0.2800$ [$-0.3600$, $-0.2000$] \\
\bottomrule
\end{tabular}
\caption{\textbf{Per-seed total compute replay on N10 ($n{=}5$ seeds, $B{=}2000$ bootstrap).}
C2 Dualformer-Auto is the only controller that saves compute vs the
always-$G{=}8$ baseline; C1 zvf-triage and C3 Hybrid both spend more.
The Dualformer-Auto CI on savings ($[+0.10, +0.18]$) and the
zvf-triage CI on savings ($[-0.36, -0.20]$) are disjoint, so the
ordering Dualformer $<$ baseline $<$ zvf-triage holds with probability
$\geq 1 - 2 \times 0.025 = 0.95$.}
\label{tab:p7-n10unif}
\end{table}

\paragraph{Paired bootstrap-CI on savings contrasts (the headline falsifiable claim).}
Treating the $5$ seeds as iid draws from a hypothetical seed population
and bootstrapping $B{=}2000$ replicate samples per contrast:

\begin{table}[h]
\centering
\small
\begin{tabular}{lrrl}
\toprule
\textbf{contrast} & \textbf{$\Delta$ mean savings} & \textbf{95\% CI} & \textbf{significance} \\
\midrule
\textbf{C2 $-$ C1} (Dualformer $-$ zvf-triage) & $+0.4200$ & [$+0.3000$, $+0.5400$] & *** \\
C3 $-$ C1 (Hybrid $-$ zvf-triage)            & $+0.0000$ & [$+0.0000$, $+0.0000$] & n.s. \\
\textbf{C3 $-$ C2} (Hybrid $-$ Dualformer)    & $-0.4200$ & [$-0.5400$, $-0.3000$] & *** \\
\bottomrule
\end{tabular}
\caption{\textbf{Paired bootstrap-CI on per-seed savings contrasts.}
Dualformer-Auto saves $42$ percentage points more than zvf-triage;
the CI excludes zero. Hybrid vs zvf-triage has zero contrast
(C3 $\equiv$ C1 on this panel, see below); Hybrid loses to Dualformer
by the same $42$pp.}
\label{tab:p7-n10unif-contrasts}
\end{table}

\paragraph{Why Hybrid $\equiv$ zvf-triage on this panel.}
The unified Hybrid (C3) has two branches: an \emph{escalation} branch
on the boundary band $z_t \in [\tau, \tau{+}\delta) = [0.7, 0.9)$ and a
\emph{de-escalation} branch on the saturation band $z_t \geq \tau{+}\delta = 0.9$.
On the N10 panel the maximum per-step ZVF observed is $0.875$ (s590
step 1), so the de-escalation branch never activates and the Hybrid
reduces to "escalate on $z_t \geq 0.7$, base otherwise" --- exactly C1.
The contrast C3 $-$ C1 = $0.0000$ with CI [$0.0000$, $0.0000$] is
literally zero because every C3 $G_t$ vector is bit-identical to the
corresponding C1 $G_t$ vector on every seed.

\paragraph{The unification license is panel-conditional.}
\textbf{On a panel that reaches the saturation band (zvf $\geq 0.9$), C3 would de-escalate those steps to $G_{\mathrm{des}}{=}4$, strictly dominating both C1 (which keeps them at $G_{\mathrm{esc}}{=}16$) and C2 (which already de-escalates them).} The N10 panel does not exercise this branch. The principled recommendation:
\begin{itemize}
\item Use \textbf{C2 Dualformer-Auto} when the per-step ZVF trajectory
stays mostly in the intermediate band $\tau \leq z_t < \tau{+}\delta$ ---
saves $14\%$ on N10 with CI $[+0.10, +0.18]$.
\item Use \textbf{C3 Hybrid} when the trajectory reaches the saturation
band $z_t \geq \tau{+}\delta$ --- strictly dominates both C1 and C2 on
that panel (falsifiable prediction: future iter with the
mega-manifest corpus will test this branch).
\item \textbf{Avoid C1 zvf-triage} on GRPO at $G_{\mathrm{base}}{=}8$
unless the trajectory never reaches the saturation band \emph{and}
the budget is unconstrained --- C1 strictly spends more compute than
the baseline on N10 (CI on savings $[-0.36, -0.20]$ entirely below zero).
\end{itemize}

\paragraph{Headroom-bad calibration check.}
Across all $5 \times 15 = 75$ fire decisions and all three controllers,
the number of steps fired on a saturated prompt ($z_t \geq 0.99$) is
exactly $0$ (CI $[0, 0]$ per controller). The trigger is well-calibrated:
it never escalates when escalation would have no value.

\paragraph{Reproduction.}
\texttt{scripts/p5p8/p7\_dualformer\_n10\_seed.py} ($\leq 280$ LoC,
stdlib only); outputs in
\texttt{experiments/results/p5p8/p7\_dualformer\_n10\_\{per\_seed.tsv,
summary.json\}}. Seed $27$, $n_{\mathrm{boot}}{=}2000$.

\subsection{Panel-conditional unification validated on the N2
saturation-band panel}
\label{sec:p7-controller-hybrid-n2}

The previous subsection closed with a falsifiable prediction: Hybrid (C3)
will strictly dominate zvf-triage (C1) only on evidence bases whose
per-step ZVF trajectory reaches the saturation band $z_t \geq \tau + \delta$.
The N10 panel does not (max $z_t = 0.875$, item~34); this subsection
replays the same three controllers on the N2 four-method tensors, where
\texttt{gift}'s trajectory \emph{does} reach the saturation band (max
$z_t = 1.000$, 8 of 40 steps at $z_t \geq 0.9$) and the falsifiable
prediction is testable for the first time.

\paragraph{Per-method compute replay on N2.}
We replay C1 (zvf-triage@$0.7$), C2 (Dualformer-Auto@$0.7$), and
C3 (Hybrid@$0.7{+}0.2$) per-(method, step) on the N2 four-method
40-step trajectories ($G_{\mathrm{base}}{=}8$, total $4 \times 40 = 160$
step-units):

\begin{table}[h]
\centering
\small
\begin{tabular}{lrrrrr}
\toprule
\textbf{method} & \textbf{zvf range} & $n_{\mathrm{sat}}$ & \textbf{C1 total} & \textbf{C2 total} & \textbf{C3 total} \\
\midrule
\texttt{grpo}  & [0.500, 0.938] & 2  & 480 & 240 & 456 \\
\texttt{aero}  & [0.562, 0.938] & 1  & 472 & 244 & 460 \\
\textbf{\texttt{gift}}  & [0.562, \textbf{1.000}] & \textbf{8}  & \textbf{528} & \textbf{216} & \textbf{432} \\
\texttt{areal} & [0.500, 0.938] & 1  & 456 & 252 & 444 \\
\midrule
\textbf{pooled (4$\times$40)} & --- & \textbf{12} & \textbf{1936} & \textbf{952} & \textbf{1792} \\
\bottomrule
\end{tabular}
\caption{\textbf{Per-method total compute replay on N2 ($4 \times 40$ step-units).}
$n_{\mathrm{sat}}$ = number of steps at $z_t \geq 0.9$. \texttt{gift}
is the only N2 method that meaningfully exercises the saturation band.
\textbf{C3 strictly reduces C1's compute by 96 rollouts on \texttt{gift}
(528$\to$432)}, the \emph{first empirical confirmation} of the iter~27
unification prediction. C2 Dualformer-Auto is uniformly the cheapest
controller on every method ($21$--$32\%$ savings vs baseline) but trades
off boundary-band precision (it de-escalates the boundary band to $G{=}4$
along with the saturation band).}
\label{tab:p7-hybrid-n2}
\end{table}

\paragraph{Headline falsifiable claim (validated).}
On the saturation-band panel (\texttt{gift}, $n_{\mathrm{sat}}{=}8$),
the per-(method, step) paired bootstrap contrast
$\texttt{C3}{-}\texttt{C1}$ on the per-step $G_t$ vector is
$\Delta_{\mathrm{total}} = -96$ rollouts (95\% bootstrap CI
$[-156, -36]$, CI excludes zero, $n_{\mathrm{boot}}{=}2000$,
seed $20260704$). This is the **first falsifiable evidence** that the
unified Hybrid is Pareto-superior to zvf-triage on a real evidence
base that exercises the saturation band. The corresponding contrasts
on \texttt{grpo} ($\Delta = -24$ $[-60, 0]$, CI barely includes zero)
and \texttt{aero}/\texttt{areal} ($\Delta = -12$ $[-36, 0]$, CI
includes zero) do not reach significance because their saturation-band
exposure is only $1$--$2$ steps; the panel-conditional unification
license holds qualitatively (C3 saves at least 8 G per sat-band step)
but the per-method bootstrap cannot certify it.

\begin{table}[h]
\centering
\small
\begin{tabular}{lrrr}
\toprule
\textbf{method} & \textbf{C3$-$C1} $\Delta_{\mathrm{total}}$ (95\% CI) & \textbf{C3$-$C2} & \textbf{C2$-$C1} \\
\midrule
\texttt{grpo}  & $-24$ $[-60, 0]$ \footnotesize n.s.    & $+216$ $[+144, +288]$ *** & $-240$ $[-312, -168]$ *** \\
\texttt{aero}  & $-12$ $[-36, 0]$ \footnotesize n.s.    & $+216$ $[+144, +288]$ *** & $-228$ $[-300, -156]$ *** \\
\textbf{\texttt{gift}}  & $\mathbf{-96}$ $[\mathbf{-156}, \mathbf{-36}]$ *** & $+216$ $[+144, +288]$ *** & $-312$ $[-372, -240]$ *** \\
\texttt{areal} & $-12$ $[-36, 0]$ \footnotesize n.s.    & $+192$ $[+120, +264]$ *** & $-204$ $[-276, -132]$ *** \\
\midrule
\textbf{pooled} & $\mathbf{-144}$ $[-228, -72]$ *** & $+840$ $[+696, +996]$ *** & $-984$ $[-1140, -840]$ *** \\
\bottomrule
\end{tabular}
\caption{\textbf{Per-(method, step) paired bootstrap contrasts on N2
($n_{\mathrm{boot}}{=}2000$, seed $20260704$).} Three stars means
the 95\% CI excludes zero. \textbf{On \texttt{gift}, C3 strictly
dominates C1 with $\Delta = -96$ and CI excludes zero, the iter~27
unification prediction is validated.}}
\label{tab:p7-hybrid-n2-contrasts}
\end{table}

\paragraph{Pooled N2 headline (160 step-units).}
Across all 4 methods $\times$ 40 steps, the Hybrid C3 saves 144 rollouts
vs C1 (95\% CI $[-228, -72]$, CI excludes zero), spends 840 rollouts more
than C2 Dualformer-Auto (CI $[+696, +996]$, excludes zero), and Dualformer
saves 984 rollouts vs C1 (CI $[-1140, -840]$, excludes zero). The
**ordering C3 $<$ C1 (compute) and C2 $<$ C3 (compute) is statistically
certified on the pooled evidence base**, but the per-method
certifiability is exactly the saturation-band step count: \texttt{gift}'s
$n_{\mathrm{sat}}{=}8$ carries the C3$-$C1 effect; the others' $n_{\mathrm{sat}}{\le}2$
are insufficient.

\paragraph{C3$\equiv$C1 collapse test, panel-conditional.}
The previous subsection proved C3$\equiv$C1 bit-for-bit on the N10 panel
(N10 max $z_t = 0.875$, de-escalation branch never fires). On N2 the
de-escalation branch fires on $1, 1, 8, 2$ steps for \texttt{aero,
areal, gift, grpo} respectively (total $12$ activation steps), so
C3 and C1 are **never** bit-identical on any N2 method. The unified
Hybrid's de-escalation branch is the active structural difference
between C3 and C1, and it activates \emph{exactly} when and only when
the per-step $z_t \geq \tau + \delta = 0.9$.

\paragraph{Headroom-bad calibration.}
All 4 controllers fire on exactly $1$ headroom-bad step (zvf$\geq 0.99$)
and exactly $12$ saturation-band steps (zvf$\geq 0.9$, distributed
$2{+}1{+}8{+}1$ across \texttt{grpo/aero/gift/areal}). The Hybrid
de-escalates the $12$ saturation-band steps to $G{=}4$ (the principled
move: no value in escalating a saturated group); \texttt{zvf-triage}
wrongly escalates them to $G{=}16$ (waste of $8$ G per step);
\texttt{Dualformer-Auto} de-escalates the saturation band \emph{and} every
boundary step (everything zvf$\geq 0.7$), which is correct for compute
economy but over-de-escalates the boundary band and is the
contrast-restoration-vs-compute-economy trade-off quantified in
Section~\ref{sec:p7-postpred}.

\paragraph{Falsifiable prediction for the next iter.}
\textbf{The unified Hybrid is now validated to be strictly
Pareto-superior to zvf-triage on saturation-band panels (\texttt{gift}:
$\Delta = -96$, $95\%$ CI $[-156, -36]$ excludes zero, $18\%$ saving).}
\textbf{Future iter}: replay the same controller set on the
mega-manifest corpus (P5, 98 cells $\times$ $G \times$ temperature
$\times$ seed) once those cells carry per-step ZVF trajectories. The
prediction is that Hybrid will strictly dominate zvf-triage on
\emph{every cell with at least one saturation-band step} and collapse
to zvf-triage on cells with no saturation-band steps.

\paragraph{Reproduction.}
\texttt{scripts/p5p8/p7\_hybrid\_unification\_n2.py} ($\leq 280$ LoC,
stdlib only); outputs in
\texttt{experiments/results/p5p8/p7\_hybrid\_n2\_\{per\_step.tsv,
per\_method.tsv, summary.json\}}. Seed $20260704$, $n_{\mathrm{boot}}{=}2000$.

